% Figure for Classical Mechanics §A4.7 (double-well potential energy diagram; after University Physics Vol. 1, Example 8.10).
\begin{tikzpicture}
  \begin{axis}[nfig, width=9cm, height=6cm, axis lines=middle, xmin=-1.2, xmax=1.25, ymin=-0.62, ymax=0.55,
      xlabel={$x$ (m)}, ylabel={$U$ (J)}, xtick={-0.707,0.707}, xticklabels={$-x_Q$,$x_Q$}, ytick={-0.5,-0.25},
      xlabel style={anchor=west}, ylabel style={anchor=south}, clip=false]
    \addplot[figB, very thick, domain=-1.06:1.06, samples=120] {2*(x^4-x^2)};
    \addplot[figR, thick, dashed, domain=-1.15:1.15] {-0.25};
    \draw[figR, line width=2.2pt, opacity=0.55] (axis cs:0.383,-0.25) -- (axis cs:0.924,-0.25);
    \draw[figR, line width=2.2pt, opacity=0.55] (axis cs:-0.924,-0.25) -- (axis cs:-0.383,-0.25);
    \fill[figR] (axis cs:0.383,-0.25) circle (1.6pt);
    \fill[figR] (axis cs:0.924,-0.25) circle (1.6pt);
    \fill[figR] (axis cs:-0.383,-0.25) circle (1.6pt);
    \fill[figR] (axis cs:-0.924,-0.25) circle (1.6pt);
    \fill[figB] (axis cs:0.707,-0.5) circle (2pt);
    \fill[figB] (axis cs:-0.707,-0.5) circle (2pt);
    \draw[figB, thick, fill=white] (axis cs:0,0) circle (2pt);
    \node[font=\scriptsize, figR, anchor=west] at (axis cs:1.12,-0.25) {$E$};
    \node[font=\scriptsize, figR, anchor=south] at (axis cs:0.383,-0.24) {$x_P$};
    \node[font=\scriptsize, figR, anchor=south west] at (axis cs:0.924,-0.24) {$x_R$};
    \node[font=\scriptsize, figB, anchor=south west] at (axis cs:0.03,0.03) {unstable};
    \node[font=\scriptsize, figB, anchor=north] at (axis cs:0.707,-0.52) {stable};
    \node[font=\scriptsize, figB, anchor=north] at (axis cs:-0.707,-0.52) {stable};
  \end{axis}
\end{tikzpicture}
